% ============================================================
% content/questions.tex —— 2026 湖南卷物理 · 试题
% 来源：2026/2026_物理真题_含解析.md（物理备课大师 Word 版转录）
% 题号保留官方编号；图形见 figures/（TikZ 重绘）与 photos/（非矢量例外）
% ============================================================

\gkpart{一、选择题}

\textbf{（一）单项选择题（本题共 6 小题，每小题 4 分，共 24 分。在每小题给出的四个选项中，只有一项是符合题目要求的。）}

\begin{question}
\gkti{1}某舰载机起飞时，在第~2~s~内的 $v$--$t$~图像如图所示，该段时间内舰载机加速度的大小为\gkblank

\begin{gkopts}(4)
  \task A．\SI{10}{\metre\per\second\squared}
  \task B．\SI{20}{\metre\per\second\squared}
  \task C．\SI{30}{\metre\per\second\squared}
  \task D．\SI{40}{\metre\per\second\squared}
\end{gkopts}

\gkfig{figures/q1-vt}
\end{question}

\begin{question}
\gkti{2}关于原子核衰变，下列说法正确的是\gkblank

\begin{gkopts}(1)
  \task A．原子核发生 $\alpha$~衰变时，电荷数增大
  \task B．原子核发生 $\alpha$~衰变时，质量数增大
  \task C．衰变产生的 $\alpha$~粒子穿透能力强，可以穿透几厘米厚的铅板
  \task D．理论研究表明，${}_{86}^{216}\mathrm{Rn}$（氡核）可能在一次衰变过程中放出两个 $\alpha$~粒子，则其衰变方程为
        ${}_{86}^{216}\mathrm{Rn} \to 2\,{}_{2}^{4}\mathrm{He} + {}_{82}^{208}\mathrm{Pb}$
\end{gkopts}
\end{question}

\begin{question}
\gkti{3}郭守敬望远镜是我国首个天文领域大科学装置，积累了大量的观测数据。分析观测数据表明，某行星绕一恒星做匀速圆周运动的周期为~$T$，轨道半径为该恒星半径的~$n$~倍。不考虑其他星体的影响，引力常量为~$G$，则该恒星的平均密度为\gkblank

\begin{gkopts}(2)
  \task A．$\dfrac{3\pi n^{3}}{GT}$
  \task B．$\dfrac{3\pi}{GT}$
  \task C．$\dfrac{3\pi n^{3}}{GT^{2}}$
  \task D．$\dfrac{3\pi}{GT^{2}}$
\end{gkopts}
\end{question}

\begin{question}
\gkti{4}如图，均匀介质中有且仅有一个点波源产生简谐横波在~$xy$~平面内传播，$A(4,\,3)$、$B(4,\,0)$~在~$xy$~平面内。某时刻，$A$~处质点位于波峰，$B$~处质点位于波谷。下列说法正确的是\gkblank

\gkfig{figures/q4-wave}

\begin{gkopts}(1)
  \task A．若波源在 $(0,\,0)$ 处，波长可能为 \SI{1}{\metre}
  \task B．若波源在 $(0,\,0)$ 处，波长可能为 \SI{2}{\metre}
  \task C．若波源在 $(0,\,3)$ 处，波长可能为 \SI{3}{\metre}
  \task D．若波源在 $(0,\,3)$ 处，波长可能为 \SI{4}{\metre}
\end{gkopts}
\end{question}

\begin{question}
\gkti{5}如图，电阻~$R$~的阻值为 \SI{10}{\ohm}，电流表为理想电表，两个电池组完全相同。若将~$a$、$b$~端分别接入~$d$、$e$~端，电流表示数为 \SI{1.0}{\ampere}；若将~$a$、$c$~端分别接入~$d$、$e$~端，电流表示数为 \SI{1.5}{\ampere}。则单个电池组的电动势和内阻分别为\gkblank

\gkfig{figures/q5-circuit}

\begin{gkopts}(2)
  \task A．\SI{15}{\volt}，\SI{5}{\ohm}
  \task B．\SI{15}{\volt}，\SI{10}{\ohm}
  \task C．\SI{30}{\volt}，\SI{5}{\ohm}
  \task D．\SI{30}{\volt}，\SI{10}{\ohm}
\end{gkopts}
\end{question}

\begin{question}
\gkti{6}如图，三个足够大的金属板~$A$、$B$、$C$~平行放置，$B$~板到~$A$、$C$~板的距离分别为~$d_{1}$、$d_{2}$，$d_{1}<d_{2}$。$a$~点到~$A$、$B$~板的距离相等，$b$~点到~$B$、$C$~板的距离相等。$A$、$C$~板均接地，$B$~板带电。关于~$a$、$b$~两点的电场强度~$E_{a}$、$E_{b}$~和电势~$\varphi_{a}$、$\varphi_{b}$，下列关系正确的是\gkblank

\gkfig{figures/q6-plates}

\begin{gkopts}(2)
  \task A．$E_{a}>E_{b}$
  \task B．$E_{a}<E_{b}$
  \task C．$\varphi_{a}>\varphi_{b}$
  \task D．$\varphi_{a}<\varphi_{b}$
\end{gkopts}
\end{question}

\textbf{（二）多项选择题（本题共 4 小题，每小题 5 分，共 20 分。在每小题给出的四个选项中，有多项符合题目要求；全部选对的得 5 分，选对但不全的得 3 分，有选错的得 0 分。）}

\begin{question}
\gkti{7}如图，一心形玩具气球内密封一定质量的理想气体和一个充有同种气体的弹性小气球，心形气球体积始终不变。在心形气球内，小气球内部气体压强大于外部气体压强，整个系统导热良好。初始时，小气球的体积为心形气球体积的一半。当温度缓慢升高时，忽略温度变化对气球材料性质的影响，下列说法正确的是\gkblank

\gkphoto[0.55\textwidth]{photos/image6-heartballoon.png}{心形气球示意图（原图保留，非矢量例外）}

\begin{gkopts}(2)
  \task A．小气球外部气体压强不变
  \task B．小气球内部气体分子数与外部相等
  \task C．小气球内部气体体积不变
  \task D．小气球内部气体体积变大
\end{gkopts}
\end{question}

\begin{question}
\gkti{8}如图，一块高功率芯片上方紧贴着一个均热板散热器。均热板是一个完全密封的扁平纯铜空腔，空腔内部注有微量的水。在正常工作过程中，水从高温芯片处吸收热量汽化，水蒸气在低温冷凝端放出热量变回液态水，并回流到底部。下列说法正确的是\gkblank

\gkfig{figures/q8-heatsink}

\begin{gkopts}(1)
  \task A．空腔内高温处所有水分子的运动速率都比低温处水分子的运动速率大
  \task B．一定量的水吸收热量变成相同温度的水蒸气，内能变大
  \task C．该均热板可以从高温物体吸热，向低温物体放热，不对外界做功
  \task D．该均热板可以从低温物体吸热，向高温物体放热，而不产生其他影响
\end{gkopts}
\end{question}

\begin{question}
\gkti{9}如图，某同学将两个相同的球形氦气球用等长的轻绳系在一起，拉着绳使气球随人一起做水平匀速运动。运动过程中，球心始终在绳的延长线上且球心连线水平，两绳所构成的平面与水平面的夹角为~$\theta$（$\theta<90^{\circ}$），两球间的弹力为~$F$。气球受到空气阻力的大小与相对空气的速率成正比，空气始终相对于地面静止，忽略气球形状的变化。该同学沿原方向运动速度缓慢增大，下列说法正确的是\gkblank

\gkfig{figures/q9-balloons}

\begin{gkopts}(2)
  \task A．$\theta$ 增大
  \task B．$\theta$ 减小
  \task C．$F$ 增大
  \task D．$F$ 减小
\end{gkopts}
\end{question}

\begin{question}
\gkti{10}某小组设计了一磁悬浮装置。如图，环形通电线圈固定在水平面上，其上方固定一半径为~$R$~的环形细管道，管道任意处磁场方向与竖直方向夹角为~$45^{\circ}$。质量为~$m$~的带正电小球在环形管道中以某一速率做匀速圆周运动，此时小球与管道间无弹力，重力加速度为~$g$。下列说法正确的是\gkblank

\gkfig{figures/q10-levitation}

\begin{gkopts}(1)
  \task A．从管道上方俯视，小球沿顺时针方向做圆周运动
  \task B．小球做圆周运动的周期为 $\pi\sqrt{\dfrac{R}{g}}$
  \task C．小球做圆周运动的半个周期内洛伦兹力的冲量大小为
        $m\sqrt{(2+\pi^{2})gR}$
  \task D．若小球的绕行方向不变，速率为其做匀速圆周运动速率的 2 倍，则小球与管道间的弹力大小为 $\sqrt{5}\,mg$
\end{gkopts}
\end{question}

\gkpart{二、实验题}

\textbf{（本题共 2 小题，共 16 分。）}

\begin{question}
\gkti{11}某同学借助视频分析软件进行``用单摆测量重力加速度''的实验，步骤如下：

\begin{enumerate}[leftmargin=2em]
  \item 准备好单摆，用支架将手机竖直放置。
  \item 打开手机录像，将单摆拉离平衡位置 $4^{\circ}$ 摆角由静止释放。
  \item 将录制的视频导入软件进行分析，得到摆球的 $v$--$t$ 图像，拟合后如图~1~所示，可知此单摆的周期~$T$~为\gkblank s。
  \item 如图~2，用刻度尺测量单摆摆长，该同学将刻度尺竖直放置，刻度尺 0 刻度线与单摆悬点对齐，读出\gkblank（填``$a$''``$b$''或``$c$''）位置的读数，该读数即为摆长~$l$。
  \item 重力加速度 $g=$\gkblank（用题中给出的字母``$T$''和``$l$''表示）。
  \item 该同学发现所得实验结果小于当地重力加速度，可能的原因是刻度尺 0 刻度线\gkblank（填``高于''或``低于''）悬点。
\end{enumerate}

\gkphoto[0.5\textwidth]{photos/image10-pendulum.png}{单摆实验装置（原图保留，非矢量例外）}
\end{question}

\begin{question}
\gkti{12}某探究小组用霍尔元件设计了一个测量微小位移的实验装置。

\gkfig{figures/q12-hall}

\begin{enumerate}[leftmargin=2em, label=(\arabic*)]
  \item 将螺旋测微器通过铜杆连接霍尔元件。霍尔元件伸入两块磁感应强度相同、同极相对放置的磁体间隙中，并处于两磁体中心竖直线上。以中心竖直线为~$z$~轴，其示意图如图~1~所示；霍尔元件上的导线与外部工作电路连接，其示意图如图~2~所示。
  \item 将电压表接在霍尔元件的\gkblank（填``$a$、$b$''或``$c$、$d$''）两端，闭合开关~$S$，测量其霍尔电压，电压表指针如图~3~所示，此时电压为\gkblank~mV。
  \item 旋转螺旋测微器的旋钮，使霍尔元件沿~$z$~轴移动至霍尔电压为 0 处，该处磁感应强度为 0，此时的螺旋测微器读数如图~4~所示，该读数为\gkblank~mm。
  \item 保持电流不变，旋转螺旋测微器的旋钮，读出霍尔元件在不同位置的霍尔电压，得到 10 组数据如下表所示。
  \item 在图~5~中描出第 4、5 组测量数据的坐标点，并作出 $z$--$U$ 图像。\gkblank
  \item 用此装置测量微小位移。取下螺旋测微器，将待测物体与铜杆连接，待测物体在~$z$~轴方向移动，其位移与霍尔元件的位移相等。某次测量中待测物体移动前电压表示数为 \SI{20}{\milli\volt}，移动后电压表示数为 \SI{60}{\milli\volt}，根据图像，此过程中物体位移的大小为\gkblank~mm（保留 2 位小数）。
\end{enumerate}

\begin{table}[ht]
\centering
\caption{霍尔电压随位移的测量数据}
\small
\begin{tabular}{ccccccccccc}
\toprule
组号 & 1 & 2 & 3 & 4 & 5 & 6 & 7 & 8 & 9 & 10\\
\midrule
$z$/mm & 11.670 & 11.960 & 12.302 & 12.633 & 12.952 & 13.270 & 13.594 & 13.930 & 14.246 & 14.585\\
$U$/mV & 9 & 16 & 25 & 35 & 42 & 48 & 57 & 67 & 74 & 81\\
\bottomrule
\end{tabular}
\end{table}

\gkfig{figures/q12-grid}

\gkphoto[0.45\textwidth]{photos/image12-meters.png}{电压表与螺旋测微器读数（原图保留，非矢量例外）}
\end{question}

\gkpart{三、解答题}

\textbf{（本题共 3 小题，共 40 分。解答应写出必要的文字说明、方程式和重要的演算步骤，只写出最后答案的不能得分。）}

\begin{question}
\gkti{13}车载摄像头需要有较大的拍摄角度。一摄像头由于结构限制，拍摄角度为 $120^{\circ}$。如图，将摄像头嵌入均匀透明介质，介质截面为矩形。只考虑该截面内光线传播情况，通过空气与介质间界面的折射，可将实际拍摄角度扩大。

\gkfig{figures/q13-camera}

\begin{enumerate}[leftmargin=2em, label=(\arabic*)]
  \item 若希望几乎贴着介质表面入射的光线 1 能够以图示路径恰好射入摄像头，即拍摄角度扩大为 $180^{\circ}$，求介质的折射率；
  \item 若介质折射率为 1.8，从侧后方向入射的光线 2 能够以图示路径折射之后发生一次全反射，然后恰好射入摄像头，求光线 2 的入射角 $\theta$ 的正弦值。
\end{enumerate}
\end{question}

\begin{question}
\gkti{14}某小组制作了一储能器，其两端电压~$U$~与其储存的电荷量~$q$~间的函数关系近似为 $U=kq$（$k$~为常量）。将该储能器接入如图所示电路，$a$、$b$、$c$~为固定的三个触点。两足够长的平行金属导轨固定于水平面上，电阻不计，间距为~$L$。导轨间有竖直向下的匀强磁场，磁感应强度为~$B$。质量为~$m$~的导体棒垂直导轨放置，接入电路的阻值为~$r$。电源电动势为~$E$，内阻不计。定值电阻阻值为~$R$。$t_{0}$~时刻开关~$S$~与~$a$~连接，直到导体棒做匀速运动，再于~$t_{1}$~时刻切换开关与~$b$~或~$c$~连接。运动过程中导体棒与导轨始终垂直且接触良好，忽略摩擦。

\gkfig{figures/q14-energy-storage}

\begin{enumerate}[leftmargin=2em, label=(\arabic*)]
  \item 求 $t_{0}$ 时刻导体棒加速度 $a_{0}$ 的大小；
  \item 若 $t_{1}$ 时刻开关与 $b$ 连接，储能器接入电路前电压为 0，当储能器电压为 $U_{0}$ 时（此时电路中电流不为 0），求导体棒速度 $v$ 的大小；
  \item 若 $t_{1}$ 时刻开关与 $c$ 连接，求从 $t_{0}$ 时刻起到导体棒静止的过程中，导体棒上产生的焦耳热 $Q$。
\end{enumerate}
\end{question}

\begin{question}
\gkti{15}如图，长为~$L$~的轻杆竖直放置，上端固定一质量为~$m$~的小球，下端连接于水平地面上某固定点，杆可绕该点无摩擦转动。小球内部安装了质量不计的智能弹射装置。受轻微扰动后，小球和杆从静止开始一起运动，当两者间弹力为 0 时，小球脱离轻杆，重力加速度为~$g$，不计空气阻力。

\gkfig{figures/q15-rod-ball}

\begin{enumerate}[leftmargin=2em, label=(\arabic*)]
  \item 求小球接触地面瞬间的速度 $v$ 的大小；
  \item 求小球接触地面瞬间的速度与水平面夹角 $\alpha$ 的正切值；
  \item 小球与地面碰撞前后，竖直方向分速度大小相等、方向相反，水平方向分速度相等。碰撞后瞬间，智能弹射装置工作，小球在极短时间内分裂成两部分，两部分速度方向均与小球分裂前瞬间的速度方向成 $\theta$ 角（$\theta$ 已知，且 $0<\theta<\alpha$）。设两部分质量之比为 $k$，弹射装置释放的能量为 $E$。
        \begin{enumerate}[leftmargin=2em, label=(\roman*)]
          \item 求 $E$ 与 $k$ 的关系；
          \item 当 $E$ 最小时，若分裂后两部分第一次落地时刻相同，求两部分第一次落地点的间距 $d$。
        \end{enumerate}
\end{enumerate}
\end{question}
